\documentclass[11pt]{article}
\usepackage[margin=.98in]{geometry}
\usepackage{amsmath,amssymb,amsthm,lmodern,microtype}
\usepackage[colorlinks=true,allcolors=blue,pdftitle={Two finiteness questions for heavenly elliptic curves}]{hyperref}
\hypersetup{colorlinks=true,linkcolor=blue,citecolor=blue,urlcolor=blue}
\newtheorem{theorem}{Theorem}
\newtheorem{lemma}[theorem]{Lemma}
\setlength{\emergencystretch}{2em}
\title{Two finiteness questions for heavenly elliptic curves\\
\large Cyclic cubic fields and descent to smaller fields}
\author{Research note}
\date{25 November 2026}
\begin{document}
\maketitle
\begin{abstract}
Both finiteness questions posed by O'Hara for cyclic cubic fields have affirmative answers. The uniform cubic prime bound needed in the argument is already unconditional. For every odd degree we also give a descent criterion which, with Dirichlet GRH, yields finiteness after excluding geometric base changes from every proper subfield. For the sextic fields with normal-closure group $A_5$, the same finiteness with the original rational-base-change exclusion holds for primes $\ell\ge17$.
\end{abstract}

For $E/K$ an elliptic curve and $\ell$ a rational prime, say that $E$ is
\emph{heavenly at $\ell$} if $K(E[\ell^\infty])/K(\mu_{\ell^\infty})$
is pro-$\ell$ and unramified away from $\ell$. Equivalently, $E$ has good
reduction away from $\ell$ and $K(E[\ell])/K(\mu_\ell)$ is an
$\ell$-extension \cite[Lemma 2.2]{ohara}. Work inside a fixed algebraic
closure of $\mathbb Q$.

\begin{theorem}[O'Hara's Questions A and B]
For cyclic cubic $K/\mathbb Q$ and primes $\ell\ge11$, there are finitely
many pairs $([E]_K,\ell)$ for which $E/K$ is heavenly at $\ell$ and its
geometric isomorphism class has no heavenly representative over $\mathbb Q$
at the same prime. The set of pairs $(\ell,K)$ occurring in this way is finite.
\end{theorem}
\begin{proof}
Rasmussen--Tamagawa \cite[Conjecture 2 and Proposition 7.5]{rt} give one
bound $N(1,3)$ valid for all cubic fields simultaneously. Thus $11\le\ell\le
N(1,3)$. For each such prime, McLeman--Rasmussen
\cite[Theorem 1.2, equivalently Theorem 3.1]{mr} with $d=3$ and
$K_0=\mathbb Q$ give finitely many geometric isomorphism classes. Their
condition $\ell-1>2d$ holds already at $11$. For each geometric class
and prime which occur after the rational-base-change exclusion,
O'Hara's Lemma~5.1 excludes $j=0,1728$, and Proposition~5.2 gives a finite
fiber of field-tagged $K$-isomorphism classes. Its hypotheses hold with
$p=3$ and $B=11$. A finite union proves Question~A. Taking the image
$([E]_K,\ell)\mapsto(\ell,K)$ proves Question~B.
\end{proof}
The horizontal bound used here depends only on the degree and dimension.
The uniformity in $K$ permits the finite union. The quantitative companion
paper uses Dirichlet GRH and gives a set of 25 candidate primes between
11 and 4,027. Which field-specific classes occur at each candidate is a
further arithmetic question.

\section*{A descent statement for general degrees}
For extensions of composite degree, there can be proper intermediate fields
other than $\mathbb Q$. Call a pair $(E/K,\ell)$ \emph{primitive over $K$}
in this note if no proper intermediate field $\mathbb Q\subseteq k\subsetneq K$
supports an elliptic curve heavenly at $\ell$ geometrically isomorphic to $E$.
In prime degree this is precisely O'Hara's exclusion.

\begin{lemma}[Descent inside a fixed division field]
Suppose $E/K$ is heavenly at a prime $\ell>\max\{2,[K:\mathbb Q]\}$, and
$j=j(E)\ne0,1728$. Put $k=\mathbb Q(j)$, choose any model $A_0/k$ with
invariant $j$, and set $M=k(A_0[\ell])$, $J=K\cap M$. There exists a
quadratic twist $A/J$ of $A_0$ (allowing the trivial twist) which is heavenly
at $\ell$ over $J$ and satisfies $A_K\simeq_K E$.
In particular, if the pair is primitive over $K$, then $K\subseteq M$.
\end{lemma}
\begin{proof}
Write $G_L$ for the absolute Galois group of $L$ and $\chi$ for the
mod-$\ell$ cyclotomic character. The mod-$\ell$ representation of $E$ on $G_K$
is upper triangular, with diagonal $\chi^{i_1},\chi^{i_2}$
\cite[Proposition 2.4]{ohara}. Since $j\ne0,1728$, there is a character
$\lambda:G_K\to\{\pm1\}$ giving the twist from $(A_0)_K$ to $E$.
In corresponding bases the mod-$\ell$ representation of $A_0$ on $G_K$ is
upper triangular, with diagonal $\lambda\chi^{i_1},\lambda\chi^{i_2}$.

The images of $G_K$ and $G_J$ on $A_0[\ell]$ coincide: restriction identifies
$\mathrm{Gal}(KM/K)$ with $\mathrm{Gal}(M/J)$, since $M/k$ is Galois.
The same basis therefore triangularizes $G_J$. Let $\phi_1,\phi_2$ be
its diagonal characters. The determinant $\chi$ factors through
$\mathrm{Gal}(M/J)$ by the Weil pairing. Thus the characters
$\phi_i\chi^{-i_i}$ on $G_J$ factor through that shared finite image and
both restrict to $\lambda$ on $G_K$. They agree and take values in
$\{\pm1\}$ everywhere. Twisting $(A_0)_J$ by this common character gives a
curve $A$ which becomes $E$ over $K$, and whose diagonal characters on
$G_J$ are the two cyclotomic powers. In particular,
$J(A[\ell])/J(\mu_\ell)$ has $\ell$-power degree.

For each $s\ge1$, put $F_s=J(A[\ell^s])$. The congruence kernels in
$\mathrm{GL}_2(\mathbb Z/\ell^s)$ are $\ell$-groups, so
$F_s/J(\mu_\ell)$ has $\ell$-power degree. At every finite prime of $J$
away from $\ell$, the inertia order of $F_s/J$ is consequently an
$\ell$-power. Since $KF_s/K$ is unramified there, multiplicativity of
ramification indices in the two towers shows that this inertia order
divides the corresponding ramification index in $K/J$. It is at most
$[K:J]<\ell$ and hence is one. The N\'eron--Ogg--Shafarevich criterion
gives good reduction of $A$ away from $\ell$. This proves descent. Primitivity
then forces $K=J$.
\end{proof}

\begin{theorem}[Odd-degree extension]
Assume Dirichlet GRH, as supplied by the Lean repository. Fix an odd integer
$n>1$ and let $K$ range over all degree-$n$ number fields. There are finitely
many field-tagged pairs $([E]_K,\ell)$ primitive over $K$ with $E$ heavenly
at a prime $\ell>2n+1$. For any odd prime $n$, the rational-base-change
exclusion alone suffices. The same assertions hold upon restricting to
Galois fields.
\end{theorem}
\begin{proof}
In the proof of Rasmussen--Tamagawa \cite[Theorem 5.3]{rt} take the fixed
base $F=\mathbb Q$. The Chebotarev application there uses the field
$\widetilde L=\mathbb Q(\mu_\ell)_{m_{\mathbb Q}}\widetilde F$; here
$\widetilde F=\mathbb Q$ and the subscript indicates a subfield of the
cyclotomic field. Its Dedekind zeta function factors into Dirichlet
$L$-functions, so the GRH input for this application is available. The choice
of the auxiliary prime is uniform in the degree-$n$ field $K$. Passage to
$K$ uses the existence of a place of odd residue degree at most $n$ over that
prime: this follows from $\sum e_vf_v=n$ odd. The rest of the cited proof
therefore gives a bound $\ell\le N(1,n)$ uniform in $K$. This specializes
the proof of their more general Dedekind-GRH theorem to its rational base.

For each remaining $\ell$, \cite[Theorem 3.1]{mr} gives finitely many $j$.
The exceptional invariants $0,1728$ are impossible in odd degree for
$\ell>3$. Indeed the CM-field containment used in
\cite[Lemma 5.1]{ohara} forces $K\mathbb Q(\sqrt{-3})$ or
$K\mathbb Q(i)$, respectively, into $K(E[\ell])$. The latter extension
over $K(\mu_\ell)$ has odd degree. Consequently the corresponding quadratic
field would lie in $K(\mu_\ell)$. It is disjoint from $\mathbb Q(\mu_\ell)$
because of ramification at $3$ or $2$, whereas
$[K(\mu_\ell):\mathbb Q(\mu_\ell)]$ divides the odd degree $n$. This argument
uses no normality of $K$. For each other $j$, primitivity and the lemma force
$K$ to lie in the fixed finite extension $k(A_0[\ell])$. This allows only
finitely many $K$. For each of them Shafarevich's finiteness theorem for
elliptic curves with good reduction away from $\ell$ gives finitely many
$K$-isomorphism classes. Taking the unions finishes the proof.
\end{proof}

\section*{An even-degree analogue for the two sextic families}
\begin{theorem}
Let $K$ range through all sextic fields with normal closure of group
$A_5$. There are finitely many pairs $([E]_K,\ell)$ with $\ell\ge17$,
$E/K$ heavenly at $\ell$, and no heavenly representative of the geometric
class over $\mathbb Q$ at $\ell$. For sextic fields of normal-closure
group $A_4$ the same finiteness holds when one excludes representatives
over every proper intermediate field. In each case the pairs $(\ell,K)$
form a finite set.
\end{theorem}
\begin{proof}
The companion paper, \emph{Uniform Rasmussen--Tamagawa bounds over two
families of sextic fields}, bounds $\ell$ by 4,027 for these fields. For
each such $\ell$, McLeman--Rasmussen's theorem applies with $d=6$, since
$\ell-1>12$. It remains to justify the exceptional invariants and the
proper-subfield condition.

Let $N$ be the normal closure, $G=\mathrm{Gal}(N/\mathbb Q)$, and
$F=\mathbb Q(\mu_\ell)$. If $j=0$ or $1728$, the CM-field containment
in O'Hara's proof puts $D=\mathbb Q(\sqrt{-3})$ or $\mathbb Q(i)$ in
$K(E[\ell])$. Over $KF$, this division field has odd degree, so
$D\subseteq KF\subseteq NF$. We rule this out directly. The group
$H=\mathrm{Gal}(NF/\mathbb Q)$ lies in $G\times\mathrm{Gal}(F/\mathbb Q)$
and surjects onto $G$. Its commutator subgroup is $G'\times\{1\}$:
the second factor is abelian and every commutator in $G$ lifts.
The largest abelian subfield of $NF$ is therefore $N^{G'}F$.
Its degree over $F$ divides three for $A_4$ and equals one for $A_5$.
Since $D$ is quadratic and ramifies at 2 or 3, while $F$ ramifies only
at $\ell\ge17$, the quadratic extension $DF/F$ cannot lie in that
abelian subfield. This excludes both exceptional invariants.

The descent lemma now bounds the fields for each $j,\ell$ under the
proper-subfield exclusion; Shafarevich bounds the curves for each field.
In the $A_5$ action the point stabilizer has order ten. Any strictly
larger proper subgroup would have index two or three, impossible by
simplicity of $A_5$ and the associated coset action. Hence $K$ has no
proper intermediate field except $\mathbb Q$. In the $A_4$ action the
order-two stabilizer lies in the Klein four subgroup, giving a cubic
intermediate field which the stated exclusion takes into account.
\end{proof}

\begin{thebibliography}{9}
\bibitem{ohara} S.~O'Hara, \emph{Heavenly elliptic curves over cubic number
fields}, September 2026, Questions A and B, Lemma 5.1 and Proposition 5.2.
\href{https://arxiv.org/abs/2609.35642v1}{arXiv:2609.35642v1}.
\bibitem{mr} C.~McLeman and C.~Rasmussen, \emph{Heavenly elliptic curves over
quadratic fields}, Res. Number Theory 12 (2026), article 79.
\href{https://doi.org/10.1007/s40993-026-00757-8}{doi:10.1007/s40993-026-00757-8}.
\bibitem{rt} C.~Rasmussen and A.~Tamagawa, \emph{Arithmetic of abelian
varieties with constrained torsion}, Trans. Amer. Math. Soc. 369 (2017),
2395--2424. \href{https://doi.org/10.1090/tran/6790}{doi:10.1090/tran/6790}.
\end{thebibliography}
\end{document}
